\documentclass[orangemode]{impaTeXPoster} 


% ==================================================
% Opções de modo da classe
% ==================================================
% A classe permite combinar múltiplos modos de estilo.
%
% Modos disponíveis:
% - simplemode : bege como cor principal
% - minimalmode: apenas branco e preto
% - darkmode   : ativa o tema escuro (pode ser combinado como outros modos)
% - pinkmode   : usa a paleta rosa
% - yellowmode : usa a paleta amarela
% - orangemode : usa a paleta laranja
% - greenmode  : usa a paleta verde
%
% Uso:
% \documentclass[pinkmode,darkmode]{impaTeXPoster}
%
% Funcionamento:
% - Modos de cor (pink, green, etc.) definem a paleta principal.
% - O darkmode ajusta o fundo e o contraste do texto.
% ==================================================


\usepackage[utf8]{inputenc}
\usepackage[T1]{fontenc}
\usepackage{lipsum}

\title{Poster Teste}
\author{Bruno Pereira}
\institute{IMPA Tech}
\date{2026}

\begin{document}

\maketitle

\begin{columns}

% ================= LEFT COLUMN =================

\column{0.48}

\begin{postercard}[green][green!20]{Introdução}
Este é um exemplo de seção usando a classe ImpaTeXPoster.

\lipsum[1]
\end{postercard}

\begin{postercard}[green][green!20][yellow]{Objetivo}
Mostrar como usar diferentes cores no ambiente.

\lipsum[2]
\end{postercard}

% ================= RIGHT COLUMN =================

\column{0.48}

\begin{postercard}{Resultados}
\lipsum[3]
\end{postercard}

\begin{postercard}[black][black!85][white][pink]{Resultados}
\lipsum[4]
\end{postercard}

\begin{postercard}{Referências}

\begin{thebibliography}{99}

\bibitem{popper}
POPPER, Karl.
\textit{A lógica da pesquisa científica}.
São Paulo: Cultrix, 2002.u

\bibitem{chalmers}
CHALMERS, A. F.
\textit{O que é ciência, afinal?}
São Paulo: Brasiliense, 1999.

\end{thebibliography}

\end{postercard}

\end{columns}

\end{document}